% \label{sec:methods} — 已在 main_zh.tex 中定义
% ============================================================================
\subsection{迁移对与偏移矩阵}
\label{sec:shiftmatrix}
我们研究六个有向跨语料库迁移对：
Chapman$\rightarrow$CPSC、Chapman$\rightarrow$PTB-XL、
CPSC$\rightarrow$Chapman、CPSC$\rightarrow$PTB-XL、
PTB-XL$\rightarrow$Chapman 和 PTB-XL$\rightarrow$CPSC，均在 4 类子空间上重新计算。L2 评估时偏移矩阵（无重训练）覆盖采样率 $500\rightarrow250\rightarrow125$~Hz；导联 $12\rightarrow6\rightarrow3\rightarrow2\rightarrow1$；噪声注入（BWL/MA/EM）信噪比 SNR $\{24,12,6,0,-6\}$~dB；增益 $\times0.5/\times2$。

\subsection{架构与随机种子}
我们使用两种架构：InceptionTime（主要架构，46{,}000--184{,}000 参数）和 ResNet1D（一维 ResNet-34，1D-ResNet-34）（次要架构，在全部 6 个迁移对上提供完整的主端点覆盖），每种架构以 5 个随机种子（42--46）进行训练。BiMamba 架构仅作为玩具实验纳入（$n=60$ 条记录，降采样分辨率），并延后至补充材料；全分辨率 BiMamba 运行在 RTX 5060（8~GB）上耗尽内存，因此 BiMamba 不进入主端点分析或架构鲁棒性声明。训练设置：AdamW + 余弦退火 + 预热（warmup）、梯度裁剪、训练期间温度冻结为 $T\equiv1$；温度仅作为事后温度缩放（Temperature Scaling, TS）应用。

\subsection{校准方法与场景}
我们将无重校准基线（None）与八种重校准方法进行比较：温度缩放（Temperature Scaling, TS）、逐类 Platt 校准（per-class Platt）、保序校准（Isotonic, OvR）、向量缩放（Vector scaling）、矩阵缩放（Matrix scaling）、Dirichlet 校准（Dirichlet calibration）、Saerens EM 先验自适应（EM）和 BBSE 先验自适应（BBSE prior adaptation）。在预注册的 13 种方法池中，None 基线加上上述八种重校准方法构成交付方法集；CORAL、TTA、MC-Dropout 和 Oracle 未在真实数据上交付（已注册的偏离，见局限性部分）。场景：\textbf{S1} 零样本迁移（源域校准拟合 $\rightarrow$ 目标域应用），为主分析场景（Sec.~\ref{sec:methods}）；S1$'$ 表示目标域无监督变体，其中 EM/BBSE 从无标签目标域预测中估计目标先验。

\subsection{预注册端点}
\label{sec:prereg}
\textbf{主要端点}（修订 A1.1，在协议 v2.1-A1 \cite{protocolv21a1} 中正式定义）：
\begin{equation}
    \Delta\text{ECE}_{\text{OOD}}
    = \text{ECE}_{\mathrm{raw}}^{\mathrm{OOD}} - \text{ECE}_{\mathrm{TS}}^{\mathrm{OOD}}
\end{equation}
其中 $\text{ECE}_{\mathrm{raw}}^{\mathrm{OOD}}$ 为目标语料库上使用原始 softmax 概率的 4 类置信度 SmoothECE，$\text{ECE}_{\mathrm{TS}}^{\mathrm{OOD}}$ 为经过源域校准拟合温度缩放后的同一指标。SmoothECE 带宽：$0.45\cdot(n/2000)^{-0.2}$。

\textbf{主要检验}（单次，单侧）：
\begin{equation}
    H_0: \Delta\text{ECE}_{\text{OOD}} = 0
    \quad\text{vs}\quad
    H_1: \Delta\text{ECE}_{\text{OOD}} > 0
\end{equation}
方向先验（已注册，经验性）：在带标签的源域校准划分上拟合的 TS，在经验上对类同分布（ID-like）数据产生很小的校准变化，而对发生偏移的目标域产生非平凡的收益 \cite{guo2017calibration,ovadia2019calibration}；这种不对称性是经验规律，而非定理。
\textbf{单侧与双侧的选择。} 单侧备择假设 $H_1: \Delta\text{ECE}_{\text{OOD}} > 0$ 的选取基于修订 A1 中注册的经验不对称性（在类同分布数据上，源域拟合温度在经验上产生很小的校准变化；在偏移目标域上，源域拟合温度存在失配，且通常保留非平凡的校准收益），\emph{而非}基于试点结果。我们强调，该方向先验是经验性的，并非定理：将 TS 应用于欠自信的目标域可能增加 ECE，且我们自己的数据中包含九个显著为负的单元格。试点结果仅触发了端点重新定位（修订 A1），而非方向选择。双侧备择假设将使有效样本量减半，并丢弃已注册的经验方向；单侧检验是预注册的选择，在多种子揭盲之前已归档于协议中。
患者级别聚类配对自助法（bootstrap），$B=10{,}000$（BCa 运行置信区间，BCa operating CI）\cite{efron1987bca}；分层合并：分层 $=$ 迁移对 $\times$ 架构。

\textbf{置信区间方法说明（BCa 与百分位法）。} 预注册的置信区间方法为 BCa（偏差校正与加速，bias-corrected and accelerated）\cite{efron1987bca}。在所报告的多种子验证中，运行置信区间方法为 $B=10{,}000$ 的 BCa，在全部 60 个实验网格上完整重跑（60/60 成功，元数据存储于 \texttt{methods.ts.meta}）。百分位自助法是 BCa 在偏差校正 $z_0\to 0$ 且加速 $a\to 0$ 时的大样本极限；对于 $B=10{,}000$ 且按患者聚类的配对聚类设计，百分位法与 BCa 的差异为 $O(n^{-1/2})$。在同时运行两种置信区间方法的 3 个代表性种子实验中（$B=10{,}000$ 百分位法重跑用于敏感性分析），BCa 与百分位法置信区间在符号和量级上均一致。来自 $B=200$ 的早期百分位置信区间结果（21/60 个实验）作为历史来源保留（5 个实验缺乏日志来源，见局限性部分）。本文全文报告的置信区间为 BCa 运行置信区间；百分位置信区间仅在两者均计算时作为敏感性交叉检验报告。

\textbf{边界条件端点}：
\begin{equation}
    \Delta\text{ECE}_{\text{ID}}
    = \text{ECE}_{\mathrm{raw}}^{\mathrm{ID}} - \text{ECE}_{\mathrm{TS}}^{\mathrm{ID}}
\end{equation}
预注册的期望是 $H_0: \Delta\text{ECE}_{\text{ID}} = 0$ \emph{不被拒绝}；鲁棒性判据为 $\geq 80\%$ 的单元格的 95\% 置信区间包含零。\emph{观测结果（BCa 运行置信区间，全网格重跑）：23/60 = 38.3\%——判据未满足，触发预注册分支 A1.4.2（百分位置信区间给出 29/60 = 48.3\%，结论相同；报告于 Sec.~\ref{sec:boundary}）。}

\textbf{次要端点}：$G$（相对于 oracle 的可恢复损失）、衰减 decay $=\Delta\text{ECE}_{\text{OOD}}-\Delta\text{ECE}_{\text{ID}}$（降级为描述性）、NLL、逐类 ECE（classwise-ECE）、Brier Murphy 分解、Cox 斜率/截距、预测/真实先验 $L_1$。

\textbf{族结构}：在 $6\text{ 个迁移对}\times 2\text{ 种架构}\times 13\text{ 个偏移水平}=156$ 项主要比较的族上采用 BH-FDR $q=0.05$；多种子验证作为鲁棒性证据（60 个种子实验）报告，而非作为确证族。

\subsection{预注册修订 A1（透明性声明）}
主要端点通过预注册修订 A1（协议 v2.0 $\rightarrow$ v2.1-A1）从 \texttt{decay} 重新定位至 $\Delta\text{ECE}_{\text{OOD}}$，该修订由一项独立的探索性试点触发，在该试点中 $\Delta\text{ECE}_{\text{ID}}$ 点估计接近零，使原始的 ID 与 OOD 对比坍缩。修订方向与经验观测到的 ID 与 OOD 不对称性一致（并非定理）；该修订作为预注册修订归档（Nosek et al.\ 2018），并附有修订前后的协议哈希。试点结果不进入主检验报告。本文所有假设均标记为\emph{探索性 + 鲁棒性验证}，而非确证性。

\textbf{预注册修订 A2（族缩减）。} 在修订 A2 下，原始的 156 项比较的 BH-FDR 族（6 个迁移对 $\times$ 2 种架构 $\times$ 13 个偏移水平）缩减为 12 个假设的确证族（6 个迁移对 $\times$ 2 种架构 $\times$ 1 个主要偏移水平），其余 13 个偏移水平被指定为探索性剂量-响应。在缩减后的族下，BH-12 给出 9/12 显著（均为 InceptionTime 方向），Bonferroni-12 给出 2/12。A2 修订归档于 \texttt{docs/PROTOCOL\_AMENDMENT\_A2\_}\allowbreak\texttt{FAMILY\_REDUCTION.md}。

\subsection{收益衰减的分解}
\label{sec:decomp}
我们使用反事实替换的加性近似，将保留的收益分解为斜率/截距/患病率贡献：
\begin{equation}
    \Delta\text{ECE} = \Delta_{\text{slope}}
    + \Delta_{\text{intercept}}
    + \Delta_{\text{prevalence}}
    + r_{\text{interaction}}
\end{equation}
其中每一项为反事实替换差分，$r_{\text{interaction}}$ 为链残差（在嵌套反事实下为零）。可识别性通过预注册阈值 $\tau$ 上的 Fisher 信息行列式进行检验；纠缠区域在误差图上报告，不设阈值。Shapley 归因以绝对值自助置信区间报告；份额可靠性标志（$|\Delta_{\text{total}}|<3/\sqrt{n}$）标记不可解释的份额。

\subsection{两层联合自助法}
温度拟合集重采样 $\rightarrow$ $T$ 分布 $\rightarrow$ 测试集重采样 $\rightarrow$ 联合 $\Delta\text{ECE}$ 置信区间（两层；诚实的非嵌套近似），依据协议 \S7。

\subsection{净临床价值（Net Clinical Value, NCV）}
\label{sec:ncv_methods}
净临床价值（Net Clinical Value, NCV）分析作为次要端点预注册，用于量化 TS 重校准在超越校准误差缩减之外的临床决策影响。对每个（样本，类别）对，我们比较校准前后的二元决策（阈值 $\tau=0.5$）：TP\_improved（原始负 $\to$ 校准正，真阳性）、TN\_improved（原始正 $\to$ 校准负，真阴性）、FP\_worsened（原始负 $\to$ 校准正，真阴性）和 FN\_worsened（原始正 $\to$ 校准负，真阳性）。NCV 定义为 $(\text{TP\_improved} + \text{TN\_improved} - \text{FP\_worsened} - \text{FN\_worsened}) / N_{\text{total}}$，其中 $N_{\text{total}} = N_{\text{samples}} \times N_{\text{classes}}$，采用与主端点相同的患者级别聚类配对自助法（$B=10{,}000$，BCa 运行置信区间）。NCV 与主端点一同作为次要端点报告，并采用 BH-FDR 校正（$q=0.05$）：显著为正的 NCV 表明 TS 在超越校准误差缩减之外改善了临床决策。NCV 在 ID 和 OOD 划分上均计算（E3 实验 CSV 中的 \texttt{ncv\_id} 和 \texttt{ncv\_ood} 列）。
